\documentclass{article}
\usepackage{/globals/de}
\begin{document}
\section{Styleguide}
\subsection{Code}
Sinnvoll viele wiederholten Codesegmente sollten Global sein.

\subsection{Commits}
Die Commitmessage ist \idr der Name des Kurses, gefolgt von einem Doppelpunkt, gefolgt von den Namen der veränderten oder neuen Themenbereiche.

\subsection{Graphiken}
Alles ist \verb|thick|, \idr.

Der Grundriss, die Basis, nutzt die standard Hintergrundsfarbe \verb|fg|. Besondere Werte oder Erkentnisse werden mit \verb|hl| hervorgehoben. Ist dies nicht genug, so wird \verb|hl2| genutzt. Hilfslinien sind \idr auch \verb|dashed|.

Für den Grundriss wird eine Größe von ungefähr $4\,\text{cm} \times 4\,\text{cm}$ angestrebt.
\subsection{Mathematik}
Unbekannte, irrelevante, oder repetitive Werte können mit \verb|\dots| ausgelassen werden.

\end{document}